\section*{अध्याय \ref{ch:g8:combustion-and-fuels} --- दहन: ईंधन, उत्पाद और ग्रीनहाउस गैसें}

\begin{solution}{exo:g8:combustion-and-fuels:1}
\ce{C + O2 -> CO2}।
\end{solution}

\begin{solution}{exo:g8:combustion-and-fuels:2}
पर्याप्त ऑक्सीजन होने पर सारा कार्बन कार्बन डाइऑक्साइड बन जाता है (पूर्ण)।
बहुत कम होने पर कुछ भाग कार्बन मोनोऑक्साइड या कालिख बन जाता है (अपूर्ण)।
\end{solution}

\begin{solution}{exo:g8:combustion-and-fuels:3}
उसका न रंग होता है, न गंध, न स्वाद: किसी को उसका पता नहीं चलता, फिर भी वह
विषैली है।
\end{solution}

\begin{solution}{exo:g8:combustion-and-fuels:4}
जीवाश्म: कोयला, प्राकृतिक गैस, पेट्रोल। नवीकरणीय: लकड़ी, बायोगैस, चुकंदर से
बना एथेनॉल।
\end{solution}

\begin{solution}{exo:g8:combustion-and-fuels:5}
\ce{C3H8 + 5O2 -> 3CO2 + 4H2O}।
\end{solution}

\begin{solution}{exo:g8:combustion-and-fuels:6}
कार्बन: 2 \ce{CO2}; हाइड्रोजन: 3 \ce{H2O}; दाईं ओर ऑक्सीजन $4 + 3 =
7$, जिसमें से 1 एथेनॉल से आती है, इसलिए 3 \ce{O2}:
\ce{C2H6O + 3O2 -> 2CO2 + 3H2O}।
\end{solution}

\begin{solution}{exo:g8:combustion-and-fuels:7}
कालिख, यानी कार्बन। \omterm{def:g4:what-a-fire-needs:burning}{दहन} अपूर्ण है: लौ को ऑक्सीजन की कमी
है।
\end{solution}

\begin{solution}{exo:g8:combustion-and-fuels:8}
1980 में लगभग 339 ppm और 2000 में लगभग 370 ppm: लगभग 30 ppm की वृद्धि।
\end{solution}

\begin{solution}{exo:g8:combustion-and-fuels:9}
लकड़ी का कार्बन पेड़ ने कुछ वर्ष या दशक पहले \omterm{def:g4:air-a-mixture-of-gases:air}{वायु} से लिया था: उसे \omterm{def:g4:what-a-fire-needs:burning}{जलाने} से
वही कार्बन डाइऑक्साइड लौटती है। कोयले का कार्बन लाखों वर्षों से ज़मीन के
नीचे जमा था: उसे \omterm{def:g4:what-a-fire-needs:burning}{जलाने} से \omterm{def:g4:air-a-mixture-of-gases:air}{वायु} में नई कार्बन डाइऑक्साइड
जुड़ती है।
\end{solution}

\begin{solution}{exo:g8:combustion-and-fuels:10}
\ce{2C + O2 -> 2CO}।
\end{solution}

\begin{solution}{exo:g8:combustion-and-fuels:11}
$(427 - 316)/316 = 111/316 \approx 0.35$: लगभग \qty{35}{\percent}।
\end{solution}

\begin{solution}{exo:g8:combustion-and-fuels:12}
शुरू में पर्याप्त ऑक्सीजन होती है: \ce{CH4 + 2O2 -> CO2 + 2H2O}। बंद कमरे की
ऑक्सीजन की खपत होते जाने पर लौ \omterm{def:g4:air-a-mixture-of-gases:air}{वायु} की कमी में \omterm{def:g4:what-a-fire-needs:burning}{जलने} लगती है:
\ce{2CH4 + 3O2 -> 2CO + 4H2O}। गंधहीन और विषैली कार्बन मोनोऑक्साइड
कमरे में जमा होती जाती है।
\end{solution}

\begin{solution}{pb:g8:combustion-and-fuels:1}
\textbf{1.} \ce{C3H8 + 5O2 -> 3CO2 + 4H2O}।

\textbf{2.} बाईं ओर: 3 C, 8 H, 10 O। दाईं ओर: 3 C, $4 \times 2 = 8$ H,
$3 \times 2 + 4 = 10$ O। संतुलित।

\textbf{3.} \ce{2C3H8 + 7O2 -> 6CO + 8H2O}।

\textbf{4.} \omterm{def:g8:combustion-and-fuels:complete}{पूर्ण दहन}: प्रोपेन के प्रति \omterm{def:g7:atoms-and-molecules:molecule}{अणु} ऑक्सीजन के 5 \omterm{def:g7:atoms-and-molecules:molecule}{अणु},
जबकि अपूर्ण में $7/2 = 3.5$।

\textbf{5.} बंद तंबू की \omterm{def:g4:air-a-mixture-of-gases:air}{वायु} में जल्दी ही ऑक्सीजन कम हो जाती है; लौ को अब
पर्याप्त ऑक्सीजन नहीं मिलती और वह अपूर्ण रूप से, पीली होकर, \omterm{def:g4:what-a-fire-needs:burning}{जलती} है।

\textbf{6.} कार्बन मोनोऑक्साइड: न रंग, न गंध (और न स्वाद)।

\textbf{7.} GHS02, लौ: अत्यंत ज्वलनशील। GHS04, गैस
सिलिंडर: दाब में रखी गई गैस।

\textbf{8.} उदाहरण के लिए: इसे केवल खुली \omterm{def:g4:air-a-mixture-of-gases:air}{हवा} में इस्तेमाल करो, कभी तंबू या
बंद कमरे में नहीं; इसे जल सकने वाली हर चीज़ से दूर रखो; कारतूस किसी भी लौ से
दूर बदलो।

\textbf{9.} $132 \div 44 = 3$ गुना भारी।

\textbf{10.} \omterm{def:g4:air-a-mixture-of-gases:air}{वायु} की ऑक्सीजन से, जो कार्बन डाइऑक्साइड (और पानी) में
पहुँच जाती है।

\textbf{11.} एक \omterm{def:g8:combustion-and-fuels:fossil}{जीवाश्म ईंधन} (यह प्राकृतिक गैस या तेल से आता है): इसे \omterm{def:g4:what-a-fire-needs:burning}{जलाने}
से \omterm{def:g4:air-a-mixture-of-gases:air}{वायु} में कार्बन डाइऑक्साइड जुड़ती है।

\textbf{12.} $450 \times 3 = \qty{1350}{g} = \qty{1.35}{kg}$ कार्बन
डाइऑक्साइड।
\end{solution}
